% Cote 119: offprints, reports, courses, CV and letters (typed and printed part): archive cote 119, site page 73 (PDF page 74).
% offprint. Draft compared: none.
% AI-transcribed by Claude Opus 5.5 (claude-opus-5-5, adaptive thinking, effort high), Message Batches API, batch r119_full (msgbatch_01DwZNm5Y7YsKBWyAN4sFs8w), 2026-11-13; not yet checked by eye.
\scanpage{119/73}
\mspage{C. R. Acad. Sc. Paris, t. 272, p. 389-392 (1\textsuperscript{er} février 1971). \quad Série A}

GÉOMÉTRIE ALGÉBRIQUE. --- $\otimes$-catégories et dualité de Tannaka.
Note (*) de M. Neantro Saavedra Rivano, transmise par M. Henri
Cartan.

Introduction. --- Soit S un schéma. Dans une Note précédente (${}^{1}$)
on a associé à un S-groupe affine et plat G une $\otimes$-catégorie ACU abélienne
Rep(G) munie d'un foncteur fibre $\omega^{\mathrm{G}} : \mathrm{Rep}(\mathrm{G}) \to \mathrm{Q\,coh}(\mathrm{S})$, et on a remar-%
qué que le couple $(\mathrm{Rep}(\mathrm{G}), \omega^{\mathrm{G}})$ permet de reconstituer le S-groupe G;
toutefois, Rep(G) ne détermine pas à elle seule G. Dans cette Note, on
s'intéresse à classifier ces $\otimes$-catégories indépendamment d'un foncteur
fibre choisi, et plus généralement à classifier les $\otimes$-catégories ACU abé-%
liennes $\Gamma(\mathrm{S}, \mathcal{O}_{\mathrm{S}})$-linéaires qui possèdent un foncteur fibre « localement »
pour la topologie fpqc. Des exemples de ces catégories apparaissent en
Géométrie algébrique, notamment dans la théorie des motifs de Grothen-%
dieck (${}^{2}$).

Pour simplifier les énoncés, on supposera que S est affine, $\mathrm{S} = \mathrm{Spec}(\mathrm{A})$.

1. Dualité de Tannaka.

1.1. Soit $C$ une catégorie abélienne A-linéaire avec des limites inductives.
Si E est un A-module, X un objet de $C$, on note $\mathrm{E} \otimes \mathrm{X}$ l'objet de $C$ qui repré-%
sente le foncteur $\mathrm{Y} \to \mathrm{Hom}_{\mathrm{A}}(\mathrm{E}, \mathrm{Hom}(\mathrm{X}, \mathrm{Y}))$. Si A$'$ est une A-algèbre (com-%
mutative, unifère), on note $C_{(\mathrm{A}')}$ la catégorie des « A$'$-modules de $C$ » : ses
objets sont les objets X de $C$ munis d'un homomorphisme de A-algèbres
$\mathrm{A}' \to \mathrm{End}(\mathrm{X})$, ses morphismes sont les morphismes de $C$ qui commutent
à l'action de A$'$. $C_{(\mathrm{A}')}$ est une catégorie abélienne A$'$-linéaire avec des limites
inductives, et on a un foncteur A-linéaire $i_{\mathrm{A}'/\mathrm{A}} : C \to C_{(\mathrm{A}')}$ défini par
$\mathrm{X} \mapsto \mathrm{A}' \otimes \mathrm{X}$ qui commute avec les limites inductives. Celui-ci a la propriété
universelle évidente pour les foncteurs A-linéaires commutant avec les
limites inductives de $C$ dans une catégorie A$'$-linéaire avec des limites
inductives.

Si $C$ possède une loi $\otimes$-ACU, $C_{(\mathrm{A}')}$ est canoniquement une $\otimes$-catégorie
ACU et $i_{\mathrm{A}'/\mathrm{A}}$ un $\otimes$-foncteur ACU. En tant que tel, il a encore une propriété
universelle évidente.

Définition 1.2. --- Une $\otimes$-catégorie ACU abélienne A-linéaire $C$ est
dite ind-tannakienne s'il existe une A-algèbre fidèlement plate A$'$ telle que

($a$) Le fonteur $i_{\mathrm{A}'/\mathrm{A}}$ soit fidèle et exact;

($b$) $C_{(\mathrm{A}')}$ soit $\otimes$-équivalente à une $\otimes$-catégorie $\mathrm{Rep}(\mathrm{G}')$, où G$'$ est un
A$'$-groupe affine et plat.

Les catégories ind-tannakiennes sur A constituent de façon naturelle
une 2-catégorie.

\mspage{\hand{73}}
